\documentclass[11pt]{article}

\usepackage[margin=1in]{geometry}
\usepackage{amsmath}
\usepackage{amssymb}
\usepackage{array}
\usepackage{hyperref}

\title{Stage 3B Physics Note: Modified-Cartoon Source Terms for the 5D Black String}
\author{GL-with-AI Project}
\date{\today}

\newcommand{\GRSP}{\texttt{GR\_SPACEDIM}}
\newcommand{\CHSP}{\texttt{CH\_SPACEDIM}}

\begin{document}
\maketitle

\section{Scope and Status}

This note is a polished companion source for the Stage 3B implementation note.
It is a derivation and validation roadmap only. No modified-cartoon source
terms are implemented in \texttt{BlackStringToy}; the current scaffold does not
evolve a five-dimensional black string.

The target reduction has gridded directions \(x,z\) and two hidden transverse
directions \(w_1,w_2\), with \(\GRSP=4\) and \(\CHSP=2\). The primary risk is
silent failure: a code may compile and run while omitting hidden-direction
algebraic terms.

The Stage 1 source inspection found that the required 5D/SO(3)
modified-cartoon CCZ4 source terms are absent from public GRChombo. Public
CCZ4 contains dimension-parameterized algebra, but not the project-specific
hidden-sphere source-term system required for \(\CHSP=2\), \(\GRSP=4\). This
is a statement about the black-string target in this project, not a defect in
public GRChombo's intended 3+1D use.

Thus Stage 3B is not merely checking whether these terms exist. It is planning
the derivation, implementation, and validation of source terms that must be
added in the project-specific layer.

\section{Reduced Geometry}

The transverse flat spatial metric can be written in spherical form as
\begin{equation}
  dx^2 + x^2 d\Omega_2^2 .
\end{equation}
In the hidden Cartesian-like cartoon representation, flat transverse space has
\begin{equation}
  \gamma_{ww}=1.
\end{equation}
The angular factor \(x^2\) is supplied by coordinate reconstruction, not by
setting \(h_{ww}\) or \(\gamma_{ww}\) to \(x^2\). The physical transverse sphere
radius is
\begin{equation}
  R = x\sqrt{\gamma_{ww}} = x\sqrt{\chi^{-1}h_{ww}} .
\end{equation}

\section{Spherical-to-Cartoon Map}

\begin{center}
\begin{tabular}{>{\raggedright}p{0.42\linewidth}>{\raggedright\arraybackslash}p{0.48\linewidth}}
\hline
Spherical object & Hidden/cartoon representation \\
\hline
\(g_{\theta\theta}=x^2\) & \(\gamma_{ww}=1\) plus external \(x^2\) radius factor \\
\(K_{\theta\theta}=x\beta^x\) & \(K_{ww}=\beta^x/x\) \\
\(\Gamma^x_{\theta\theta}=-x\) & hidden algebraic source terms after contraction \\
\(\Gamma^\theta_{x\theta}=1/x\) & explicit \(1/x\) reduced source terms \\
angular trace contribution & two hidden contributions, one for each suppressed direction \\
diagonal trace & \(T^x{}_x + T^z{}_z + 2T^w{}_w\) \\
\hline
\end{tabular}
\end{center}

For this project the number of suppressed spatial directions is
\begin{equation}
  N_{\rm hidden}=\GRSP-\CHSP=4-2=2,
\end{equation}
so diagonal traces contain \(2T^w{}_w\). Writing the factor this way helps
avoid hard-coding the wrong multiplicity in future variants. The two hidden
directions are equivalent here because of the imposed SO(3) spherical symmetry.
In source-code notation, this is
\verb|N_hidden = GR_SPACEDIM - CH_SPACEDIM = 4 - 2 = 2|.

\section{GP Extrinsic-Curvature Check}

For the natural GP shift
\begin{equation}
  \beta^x = s\sqrt{\frac{r_0}{x}},
\end{equation}
and flat spatial metric
\begin{equation}
  dl^2 = dx^2 + x^2 d\Omega_2^2 + dz^2,
\end{equation}
the relevant spherical Christoffels are
\begin{align}
  \Gamma^x_{\theta\theta} &= -x, &
  \Gamma^x_{\phi\phi} &= -x\sin^2\theta,\\
  \Gamma^\theta_{x\theta} &= \frac{1}{x}, &
  \Gamma^\phi_{x\phi} &= \frac{1}{x}.
\end{align}
Using
\begin{equation}
  K_{ij}
  =
  \frac{1}{2\alpha}
  \left(D_i\beta_j+D_j\beta_i-\partial_t\gamma_{ij}\right),
\end{equation}
with \(\alpha=1\), one obtains
\begin{align}
  K_{xx} &= \partial_x\beta_x
  = -\frac{s}{2}\sqrt{\frac{r_0}{x^3}},\\
  K_{\theta\theta} &= x\beta^x = s\sqrt{r_0x},\\
  K_{\phi\phi} &= x\beta^x\sin^2\theta,\\
  K_{zz} &= 0.
\end{align}
The \(\sin^2\theta\) factor cancels in the contracted contribution:
\begin{align}
  \gamma^{\phi\phi}K_{\phi\phi}
  &=
  \frac{1}{x^2\sin^2\theta}
  \left(x\beta^x\sin^2\theta\right)
  =
  \frac{\beta^x}{x},\\
  \gamma^{\theta\theta}K_{\theta\theta}
  &=
  \frac{1}{x^2}\left(x\beta^x\right)
  =
  \frac{\beta^x}{x}.
\end{align}
This demonstrates why the two suppressed directions contribute equally after
contraction and why the cartoon representation uses two copies of
\(K^w{}_w\).
Mapping to hidden cartoon components gives
\begin{equation}
  K_{ww} = \frac{\beta^x}{x} = s\sqrt{\frac{r_0}{x^3}},
\end{equation}
for each hidden direction, and
\begin{equation}
  K
  =
  K^x{}_x+K^z{}_z+K^{w_1}{}_{w_1}+K^{w_2}{}_{w_2}
  =
  \frac{3s}{2}\sqrt{\frac{r_0}{x^3}}.
\end{equation}
This check assumes the unperturbed \(\chi=1\) GP slice and remains subject to
GRChombo's possible overall \(K_{ij}\) sign convention.

\section{Source-Term Categories}

The eventual CCZ4 implementation must handle hidden metric variables
\(h_{ww}\) and \(\gamma_{ww}\), hidden extrinsic curvature \(A_{ww}\) and
\(K_{ww}\), full-dimensional traces with \(\GRSP=4\), algebraic \(1/x\) and
\(1/x^2\) replacements for suppressed derivatives, hidden Ricci tensor/scalar
terms, conformal connection functions and Gamma constraints, gauge RHS terms,
and regularization near small \(x\).

In this Stage 3B target, hidden-direction multiplicities should be written as
\(N_{\rm hidden}=\GRSP-\CHSP=2\) equivalent SO(3) contributions, not inferred
from the number of grid axes alone.

Detailed CCZ4 RHS formulas are future work and require symbolic verification
before implementation.

\section{Validation Strategy}

Validation should proceed in layers:
\begin{enumerate}
  \item Symbolically derive reduced Christoffels, Ricci terms, traces, and CCZ4
    RHS contributions, preferably with xAct/Mathematica or SymPy.
  \item Unit-test Minkowski cartoon variables, the uniform GP black string, and
    the perturbed GL seed.
  \item Track runtime diagnostics: Hamiltonian and momentum constraints,
    \(\Theta\), tracelessness, extrema of \(h_{ww}\) and \(\gamma_{ww}\),
    source-term norms, regularity near small \(x\), and horizon radius/area.
  \item Run at least three resolutions and compare to exact/unperturbed
    solutions, because wrong source terms can converge to the wrong equations.
\end{enumerate}

\section{Failure Modes}

Common silent failures include using \(\CHSP=2\) in physical traces, missing
the factor \(N_{\rm hidden}=\GRSP-\CHSP=4-2=2\) for the two equivalent hidden
directions, setting \(h_{ww}=x^2\), confusing \(K_{ww}\) with
\(K_{\theta\theta}\), dropping \(1/x\) or \(1/x^2\) terms, failing to
regularize near small \(x\), or using an enum layout that causes AH or
source-term code to read the wrong variable.

\section{Unresolved Before Coding}

Before implementation, the project still needs a staged CCZ4 modified-cartoon
RHS derivation. The future sub-plan is to derive reduced metric and
inverse-metric bookkeeping including \(h_{ww}\) and \(\gamma_{ww}\), derive
reduced Christoffels including all \(1/x\) terms and regular limits, derive
hidden Ricci tensor/scalar contributions, derive conformal-metric RHS terms
involving \(h_{ww}\), derive extrinsic-curvature RHS terms involving \(A_{ww}\)
and trace terms, derive Gamma/connection and constraint-propagation terms,
derive gauge-equation terms and verify whether hidden-direction contributions
enter, and validate each block against symbolic calculations and exact test
data before coding the next block.

The project also still needs symbolic verification, a small-\(x\)
regularization decision, the exact \texttt{hww}/\texttt{Aww} enum layout, a
reliable \(\GRSP=4\) build path, unit-test definitions, direct comparison to
spherical-coordinate formulas, and convergence acceptance criteria.

\end{document}
